\subsection{庞加莱群胚}

定义 1. --- 设 X 是一个拓扑空间，\(c_{0}\) 和 \(c_{1}\) 是 X 中的道路，\(\sigma\colon\mathbf{I}\times\mathbf{I}\to\mathbf{X}\) 是连接 \(c_{0}\) 与 \(c_{1}\) 的同伦。称 \(\sigma\) 为严格同伦，若映射 \(s\mapsto\sigma(0,s)\) 和 \(s\mapsto\sigma(1,s)\) 为常值。

如果存在连接 \(c_{0}\) 与 \(c_{1}\) 的严格同伦，就称 X 中的道路 \(c_{0}\) 和 \(c_{1}\) 严格同伦。

两条严格同伦的道路具有相同的起点和终点。

例. --- 设 c 是 X 中的一条道路，\(\varphi\colon I\to I\) 是连续映射，满足 \(\varphi(0)=0\) 且 \(\varphi(1)=1\)。道路 c 与 \(c\circ\varphi\) 严格同伦。事实上，由 \(\sigma\colon I\times I\to X\) 和 \(\sigma(t,s)=c((1-s)t+s\varphi(t))\) 定义的映射是连接 c 与 \(c\circ\varphi\) 的严格同伦。

设 X 是一个拓扑空间。回忆（参见 III, p.~257），\(\Lambda(X)\) 表示道路空间 \(\mathcal{C}_c(\mathbf{I};\mathrm{X})\)，并且对于 \(x\)、\(y \in X\)，\(\Lambda_{x,y}(X)\) 是 \(\Lambda(X)\) 的子空间，由起点为 x、终点为 y 的道路组成。对于 \(x\)、\(y\)，集合 \(\Lambda_{x,y}(X)\) 是 X 的道路空间的一个划分。根据（III, p.~257, 注 2）从 \(x\)、\(y \in X\)、\(\mathcal{C}(\mathbf{I} \times \mathbf{I};\mathrm{X})\) 到 \(\mathcal{C}(\mathbf{I};\Lambda(\mathrm{X}))\) 的典范双射，严格同伦对应于这样的道路 \(c: \mathbf{I} \to \Lambda(X)\)，其像包含于形如 \(\Lambda_{x,y}(\mathrm{X})\) 的子空间中。因此，关系“道路 \(c_0\) 和 \(c_1\) 严格同伦”是 \(\Lambda(\mathrm{X})\) 上的等价关系（III, p.~259, 命题 5），而该关系的等价类是 X 的道路空间的子空间 \(\Lambda_{x,y}(\mathrm{X})\) 的道路连通分支。记 \(\varpi_{x,y}(\mathrm{X})\) 为集合 \(\pi_0(\Lambda_{x,y}(\mathrm{X}))\)，并称连接 \(x\) 与 y 的道路类为 \(y\)、\(\varpi_{x,y}(\mathrm{X})\) 的任一元素。

定义 2. --- 我们称 X 的圈空间为 \(\Omega(X)\)，它是 \(\Lambda(X)\) 的由 X 中的圈组成的子空间（III, p.~256, 定义 1）。

记 \(\Omega_x(\mathrm{X})\) 为集合 \(\Lambda_{x,x}(\mathrm{X})\)。\(\Omega_x(\mathrm{X})\) 的元素称为 X 中 \(x\) 处的圈，\(\varpi_{x,x}(\mathrm{X})\) 的元素称为 X 中 \(x\) 处的圈类。映射 \(e_x: \mathbf{I} \to \mathbf{X}\) 以 \(x\) 为像中的唯一点，是一个圈，称为 \(x\) 处的常值圈；其严格同伦类记为 \(\varepsilon_x\)。映射 \(x \mapsto e_x\) 从 X 到 \(\Lambda(\mathrm{X})\) 是连续的（III, p.~257, 命题 1）。

设 X 是一个拓扑空间，x、y、z 是 X 中的点。由连续映射 \(c \mapsto \overline{c}\)（从 \(\Lambda_{x,y}(X)\) 到 \(\Lambda_{y,x}(X)\)）（III, p.~258, 推论）取道路连通分支，得到从 \(\varpi_{x,y}(X)\) 到 \(\varpi_{y,x}(X)\) 的映射，记为 \(\gamma \mapsto \overline{\gamma}\)。若 \(\gamma \in \varpi_{x,y}(X)\)，则称 \(\overline{\gamma}\) 为道路类 \(\gamma\) 的逆。

同样，在将集合 \(\pi_0(\Lambda_{x,y}(\mathrm{X})) \times \pi_0(\Lambda_{y,z}(\mathrm{X}))\) 与 \(\pi_0(\Lambda_{x,y}(\mathrm{X}) \times \Lambda_{y,z}(\mathrm{X}))\) 识别之后（III, p.~260, 命题 6），由连续映射 \((c,d) \mapsto c * d\) 从 \(\Lambda_{x,y}(\mathrm{X}) \times \Lambda_{y,z}(\mathrm{X})\) 到 \(\Lambda_{x,z}(\mathrm{X})\)（III, p.~258, 推论）取道路连通分支后得到映射 \(C_{x,y,z} \colon \varpi_{x,y}(\mathrm{X}) \times \varpi_{y,z}(\mathrm{X}) \to \varpi_{x,z}(\mathrm{X})\)。对于 \(\gamma \in \varpi_{x,y}(\mathrm{X})\) 和 \(\delta \in \varpi_{y,z}(\mathrm{X})\)，记 \(\gamma \delta\) 为严格同伦类 \(C_{x,y,z}(\gamma,\delta)\)。称其为可拼接道路类 \(\gamma\) 与 \(\delta\) 的复合。

有 \(\overline{\overline{\gamma}}=\gamma\) 以及 \(\overline{\gamma\delta}=\overline{\delta}\overline{\gamma}\)。

命题 1. --- 设 X 是一个拓扑空间，x、y、z、u 是 X 中的点，\(\gamma_{1} \in \varpi_{x,y}(X)\)、\(\gamma_{2} \in \varpi_{y,z}(X)\)、\(\gamma_{3} \in \varpi_{z,u}(X)\) 是道路类。有

\begin{enumerate}
\def\labelenumi{(\arabic{enumi})}
\tightlist
\item
\end{enumerate}

\[
\varepsilon_ {x} \gamma_ {1} = \gamma_ {1} \varepsilon_ {y} = \gamma_ {1},\tag{2}
\]

\[
\gamma_ {1} \overline {{\gamma}} _ {1} = \varepsilon_ {x}, \qquad \overline {{\gamma}} _ {1} \gamma_ {1} = \varepsilon_ {y},\tag{3}
\]

\[
(\gamma_ {1} \gamma_ {2}) \gamma_ {3} = \gamma_ {1} (\gamma_ {2} \gamma_ {3}).
\]

令 \(c_{1} \in \Lambda_{x,y}(\mathrm{X})\)、\(c_{2} \in \Lambda_{y,z}(\mathrm{X})\)、\(c_{3} \in \Lambda_{z,u}(\mathrm{X})\) 分别为 \(\gamma_{1}\)、\(\gamma_{2}\) 和 \(\gamma_{3}\) 的代表元。令 \(\varphi: \mathbf{I} \to \mathbf{I}\) 为下式定义的函数：

\[
\varphi (t) = \left\{ \begin{array}{l l} t / 2 & \text { for } 0 \leqslant t \leqslant 1 / 2, \\ t - 1 / 4 & \text { for } 1 / 2 \leqslant t \leqslant 3 / 4, \\ 2 t - 1 & \text { for } 3 / 4 \leqslant t \leqslant 1. \end{array} \right.\tag{4}
\]

函数 \(\varphi\) 在三个区间 \([0,1/2]\)、\([1/2,3/4]\) 和 \([3/4,1]\) 上分别是仿射的，因而连续。有 \(\varphi(0) = 0\) 和 \(\varphi(1) = 1\)。由 III, p.~256 定义道路拼接的公式 (1)，可得

\[
c _ {1} * (c _ {2} * c _ {3}) = ((c _ {1} * c _ {2}) * c _ {3}) \circ \varphi .
\]

根据 III, p.~289 的例，路径 \(c_{1} * (c_{2} * c_{3})\) 与 \((c_{1} * c_{2}) * c_{3}\) 严格同伦，故等式 (3) 成立。

同样，由下式定义的函数 \(\psi\colon\mathbf{I}\to\mathbf{I}\)

\[
\psi (t) = \left\{ \begin{array}{l l} 2 t & \text { for } 0 \leqslant t \leqslant 1 / 2, \\ 1 & \text { for } 1 / 2 \leqslant t \leqslant 1 \end{array} \right.\tag{5}
\]

是连续的，并满足 \(\psi(0)=0,\psi(1)=1\)。等式 \(\gamma_{1}\varepsilon_{y}=\gamma_{1}\) 由下述事实得到：

\[
c _ {1} * e _ {y} = c _ {1} \circ \psi
\]

而等式 \(\varepsilon_{x}\gamma_{1}=\gamma_{1}\) 的证明相同，故得 (1)。

由下式定义的映射 \(\sigma\colon\mathbf{I}\times\mathbf{I}\to\mathrm{X}\)

\[
\sigma (t, s) = \left\{ \begin{array}{l l} c _ {1} (2 t s) & \text { for } 0 \leqslant t \leqslant 1 / 2, \\ c _ {1} (2 (1 - t) s) & \text { for } 1 / 2 \leqslant t \leqslant 1 \end{array} \right.\tag{6}
\]

是连续的；它是连接道路 \(e_{x}\) 与道路 \(c_{1}*\overline{c_{1}}\) 的严格同伦，故得到 (2) 中的第一个等式。第二个等式由第一个等式以及对每条道路 c 都有 \(c=\overline{\overline{c}}\) 这一事实得到。

注 1. --- 设 X 是一个拓扑空间。设 n 是满足 \(\geqslant 1\) 的整数，\((c_{1},\ldots,c_{n})\) 是 X 中的道路序列，对 \(c_{i}\)，\(c_{i+1}\) 与 \(1 \leqslant i \leqslant n-1\) 可拼接（这样的序列称为可拼接道路序列）。记道路 c 为

\[
c _ {1} * \left(c _ {2} * \left(\dots * \left(c _ {n - 1} * c _ {n}\right) \ldots\right)\right)
\]

以及由 \(c'\)（其中 \(c'(t) = c_i(nt - i + 1)\) 且 \(1 \leqslant i \leqslant n\)）定义的道路 \(t \in [\frac{i-1}{n}, \frac{i}{n}]\)。道路 \(c\) 与 \(c'\) 具有相同的像并且严格同伦：其中一条是另一条与一个固定 0 和 1 的 I 的同胚的复合（参见 III, p. 289, 例）。我们有时将道路 \(c_{1} * c_{2} * \cdots * c_{n}\) 记为 \(c'\)。

存在唯一的有向图 \(\varpi(\mathrm{X})\)，其顶点集为 X，连接点 \(x\) 与点 \(y\) 的箭集为 \(\varpi_{x,y}(\mathrm{X})\)，且映射 \(C_{x,y,z}\) 定义一个合成法则。由命题 1，\(\varpi(\mathrm{X})\) 是一个群胚（II, p.~162, 定义 4）。对每个 \(x \in \mathrm{X}\)，该群胚在顶点 x 处的单位元是像为 x 的常值圈的类。箭 \(x\)、\(x\)、\(\gamma\) 的逆是箭 \(\overline{\gamma}\)，我们也记为 \(\gamma^{-1}\)。特别地，合成法则 \(C_{x,x,x}\) 赋予每个 \(x \in \mathrm{X}\) 的集合 \(\varpi_{x,x}(\mathrm{X})\) 一个群结构；将此群记为 \(\pi_1(\mathrm{X}, x)\)。

定义 3. --- 设 X 是一个拓扑空间。群胚 \(\varpi(X)\) 称为 X 的庞加莱群胚或基本群胚。设 x 是 X 中的一点；x 处圈类组成的群 \(\pi_1(X,x)\) 称为 X 在点 x 处的庞加莱群或基本群。

设 \(\mathcal{U}\) 是 X 的一个子集族，其内部覆盖 X。X 中像包含于 \(\mathcal{U}\) 的某个成员中的道路类生成群胚 \(\varpi(\mathrm{X})\)（III, p.~272 的引理 4）。

群胚 \(\varpi(X)\) 的轨道（II, p.~162）与空间 X 的道路连通分支一致。特别地（上引文），有：

命题 2. --- 设 X 是一个拓扑空间。

\begin{enumerate}
\def\labelenumi{\alph{enumi})}
\item
  如果空间 X 是道路连通的，则对 \(\pi_1(\mathrm{X},x)\) 中的所有点 \(\pi_1(\mathrm{X},y)\) 和 \(x\)，群 \(y\) 与 \(\mathrm{X}\) 同构。
\item
  设 x 是 X 中的一点；下列条件等价：
\end{enumerate}

\begin{enumerate}
\def\labelenumi{(\roman{enumi})}
\item
  群 \(\pi_1(X,x)\) 是平凡的；
\item
  X 中起点为 x 且具有相同终点的两条道路严格同伦；
\item
  X 中每个起点为 x 的圈都与像为 x 的常值圈严格同伦。
\end{enumerate}

更确切地说，设 X 是一个道路连通的拓扑空间，x 和 y 是 X 中的点。对 \(\delta\) 的每个元素 \(\varpi_{x,y}(X)\)，映射 \(u_{\delta} \colon \gamma \mapsto \delta \gamma \delta^{-1}\) 从 \(\pi_1(\mathrm{X},y)\) 到 \(\pi_1(\mathrm{X},x)\) 是群同构，其逆同构为 \(u_{\delta^{-1}}\)。对于 \(\delta \in \varpi_{x,y}(\mathrm{X})\) 和 \(\gamma \in \pi_1(\mathrm{X},y)\)，置 \(^{\delta}\gamma = u_{\delta}(\gamma)\)。当 \(x = y\) 时，有 \(^{\delta}\gamma = \operatorname{Int}(\delta)(\gamma)\)，即 \(u_{\delta} = \operatorname{Int}(\delta)\)。

设 \(x, y, z\) 是 X 中的点。对于 \(\delta \in \varpi_{x,y}(X)\) 和 \(\eta \in \varpi_{y,z}(X)\)，有 \(u_{\delta \eta} = u_{\delta} \circ u_{\eta}\)，这也可对 \(\delta \eta \gamma = {}^{\delta}(\eta \gamma)\) 写成 \(\gamma \in \pi_1(X, z)\)。

设 \(x, y\) 是 X 中的点，\(\delta\)、\(\delta'\) 是 \(\varpi_{x,y}(\mathrm{X})\) 的元素。有

\[
u _ {\delta^ {\prime}} = u _ {\delta} \circ \operatorname{Int} (\delta^ {- 1} \delta^ {\prime}) = \operatorname{Int} (\delta^ {\prime} \delta^ {- 1}) \circ u _ {\delta}.\tag{7}
\]

注. --- 2) 设 X 是一个拓扑空间，C 是 X 的一个道路连通分支。若 x 和 y 是 C 中的点，则拓扑空间 \(\Lambda_{x,y}(\mathrm{C})\) 可与拓扑空间 \(\Lambda_{x,y}(\mathrm{X})\) 识别，从而集合 \(\varpi_{x,y}(\mathrm{C})\) 可与集合 \(\varpi_{x,y}(\mathrm{X})\) 识别。因此，基本群胚 \(\varpi(\mathrm{C})\) 可与 \(\varpi(\mathrm{X})\) 的、以 C 为点集的全子群胚识别。特别地，对 C 中的每个点 x，群 \(\pi_1(\mathrm{C}, x)\) 可与群 \(\pi_1(\mathrm{X}, x)\) 识别。

\begin{enumerate}
\def\labelenumi{\arabic{enumi})}
\setcounter{enumi}{2}
\tightlist
\item
  设 X 是一个拓扑空间，\(x\)、\(y\)、\(z\) 是 X 中的点。从 \((c,d)\mapsto c*d\) 到 \(\Lambda_{x,y}(\mathrm{X})\times \Lambda_{y,z}(\mathrm{X})\) 的映射 \(\Lambda_{x,z}(\mathrm{X})\) 是连续的（III, p.~258, 命题 2 的推论）。注意，当赋予集合 \(\varpi_{x,y}(\mathrm{X})\times \varpi_{y,z}(\mathrm{X})\to \varpi_{x,z}(\mathrm{X})\)、\(\varpi_{x,y}(\mathrm{X})\) 和 \(\varpi_{y,z}(\mathrm{X})\) 商拓扑时，由它诱导的合成映射 \(\varpi_{x,z}(\mathrm{X})\) 不一定连续（参见 TG, I, p.~35 及 III, p.~259）。然而，对每个 \(\gamma_0\in \varpi_{x,y}(\mathrm{X})\) 和每个 \(\delta_0\in \varpi_{y,z}(\mathrm{X})\)，部分映射 \(\gamma \mapsto \gamma \delta_0\) 从 \(\varpi_{x,y}(\mathrm{X})\) 到 \(\varpi_{x,z}(\mathrm{X})\)，以及 \(\delta \mapsto \gamma_0\delta\) 从 \(\varpi_{y,z}(\mathrm{X})\) 到 \(\varpi_{x,z}(\mathrm{X})\)，都是同胚。事实上，令 \(c_0\in \Lambda_{x,y}(\mathrm{X})\) 为类为 \(\gamma_0\) 的道路。映射 \(d\mapsto c_0*d\) 是从 \(\Lambda_{y,z}(\mathrm{X})\) 到 \(\Lambda_{x,z}(\mathrm{X})\) 的连续映射。映射 \(\delta \mapsto \gamma_0\delta\) 通过取商得到，因而连续。从 \(\delta^{\prime}\mapsto \gamma_{0}^{-1}\delta^{\prime}\) 到 \(\varpi_{x,z}(\mathrm{X})\) 的映射 \(\varpi_{y,z}(\mathrm{X})\) 同样如此。由于这两个映射互为逆映射，它们是同胚。映射 \(\gamma \mapsto \gamma \delta_0\) 的情形类似。另见 IV, p.~374。
\end{enumerate}

